% --- Archon physics formalization source begin ---
% archon:physics
% archon:covers PhyXMiniProblems/problem_phyx_mini_0191.lean
% archon:source-report reports/phyx_mini/problem_phyx_mini_0191.source.json

\chapter{Physics problem phyx\_mini\_0191}
\label{ch:PhyXMiniProblems_problem_phyx_mini_0191}

\paragraph{Problem source.}
\#\# Physical scenario

Consider an idealized bird (treated as a point source) that emits constant sound power, with intensity obeying the inverse-square law (figure).You move twice the distance from the bird.

\#\# Question

How many decibels does the sound intensity level drop?

\#\# Answer choices

- A: A: 6.0dB

- B: B: 8.2dB

- C: C: 7.5dB

- D: D: 9.1dB

\#\# Figure caption (auxiliary; use the image as primary evidence)

The image illustrates a scene with two human figures labeled as \textbackslash{}( P\_1 \textbackslash{}) and \textbackslash{}( P\_2 \textbackslash{}), a bird, and a series of concentric blue circles. The bird is depicted as the "Point source" of sound waves, represented by the circles emanating from it. The bird is perched on a rock and shown singing with musical notes above its head.

The human figure labeled \textbackslash{}( P\_1 \textbackslash{}) is closer to the bird, while \textbackslash{}( P\_2 \textbackslash{}) is further away. The circles indicate sound waves spreading outwards from the bird. The background features a simple landscape with trees or shrubs. The overall scene visually represents the concept of sound waves radiating from a point source.

\#\# Dataset metadata

Sound; Physical Model Grounding Reasoning, Multi-Formula Reasoning

\paragraph{Recorded answer/context.}
A: A: 6.0dB

\paragraph{Figure/image path.}
/root/proposal\_for\_physic/science-mango/phyx\_mini\_run/phyx\_data/test\_image/191.png

\paragraph{Formalization target.}
create a compiling Lean file with sorry bodies at `PhyXMiniProblems/problem\_phyx\_mini\_0191.lean`. The Lean declarations must preserve the physical quantities, dimensions or dimensional roles, figure labels, governing-law hypotheses, and final relation expressed by this problem.
Use Mathlib/Physlib names found through LeanExplore where available. If a physics API is missing, introduce faithful local abstractions rather than scalar placeholder aliases.

\begin{theorem}[Physics formalization target]
\label{thm:physics:phyx_mini_0191:target}
\lean{PhyXMiniProblems.ProblemPhyXMini0191.problem_phyx_mini_0191}
\uses{def:physics:phyx-mini-0191:phyxminiproblems-problemphyxmini0191-birdpointsourcesoundsetup, def:physics:phyx-mini-0191:phyxminiproblems-problemphyxmini0191-matchespointsourcescenario, def:physics:phyx-mini-0191:phyxminiproblems-problemphyxmini0191-matchessuppliedfigure, def:physics:phyx-mini-0191:phyxminiproblems-problemphyxmini0191-movedtotwicethedistance, def:physics:phyx-mini-0191:phyxminiproblems-problemphyxmini0191-hasphysicalacousticparameters, def:physics:phyx-mini-0191:phyxminiproblems-problemphyxmini0191-satisfiesisotropicinversesquarelaw, def:physics:phyx-mini-0191:phyxminiproblems-problemphyxmini0191-soundleveldropindecibels, def:physics:phyx-mini-0191:phyxminiproblems-problemphyxmini0191-matchesanswertonearesttenthdecibel}
The assigned autoformalize agent should translate this physics problem into Lean declarations in the covered file, with theorem and lemma proof bodies written as `by sorry`.
\end{theorem}
\begin{proof}
This is an autoformalization task, not a proof task. Produce statements that can later be proved by the physics prover without weakening the physical model.
\end{proof}

% --- Archon declaration topology begin ---
\section{Declaration topology}
\begin{definition}[Length Magnitude]
  \label{def:physics:phyx-mini-0191:phyxminiproblems-problemphyxmini0191-lengthmagnitude}
  \lean{PhyXMiniProblems.ProblemPhyXMini0191.LengthMagnitude}
  A nonnegative physical distance from the point source
\end{definition}
\begin{proof}
  This is the declared model definition of Length Magnitude.
\end{proof}
\begin{definition}[Acoustic Power Quantity]
  \label{def:physics:phyx-mini-0191:phyxminiproblems-problemphyxmini0191-acousticpowerquantity}
  \lean{PhyXMiniProblems.ProblemPhyXMini0191.AcousticPowerQuantity}
  Acoustic power has SI unit watt and dimension mass * length\textasciicircum{}2 / time\textasciicircum{}3
\end{definition}
\begin{proof}
  This is the declared model definition of Acoustic Power Quantity.
\end{proof}
\begin{definition}[Acoustic Intensity Quantity]
  \label{def:physics:phyx-mini-0191:phyxminiproblems-problemphyxmini0191-acousticintensityquantity}
  \lean{PhyXMiniProblems.ProblemPhyXMini0191.AcousticIntensityQuantity}
  Acoustic intensity is power per area, with SI unit watt per square metre and dimension mass / time\textasciicircum{}3
\end{definition}
\begin{proof}
  This is the declared model definition of Acoustic Intensity Quantity.
\end{proof}
\begin{definition}[distance Readout]
  \label{def:physics:phyx-mini-0191:phyxminiproblems-problemphyxmini0191-distancereadout}
  \lean{PhyXMiniProblems.ProblemPhyXMini0191.distanceReadout}
  \uses{def:physics:phyx-mini-0191:phyxminiproblems-problemphyxmini0191-lengthmagnitude}
  Read a physical distance in a coherent choice of units
\end{definition}
\begin{proof}
  This is the declared model definition of distance Readout.
\end{proof}
\begin{definition}[acoustic Power Readout]
  \label{def:physics:phyx-mini-0191:phyxminiproblems-problemphyxmini0191-acousticpowerreadout}
  \lean{PhyXMiniProblems.ProblemPhyXMini0191.acousticPowerReadout}
  \uses{def:physics:phyx-mini-0191:phyxminiproblems-problemphyxmini0191-acousticpowerquantity}
  Read an acoustic power in a coherent choice of units
\end{definition}
\begin{proof}
  This is the declared model definition of acoustic Power Readout.
\end{proof}
\begin{definition}[acoustic Intensity Readout]
  \label{def:physics:phyx-mini-0191:phyxminiproblems-problemphyxmini0191-acousticintensityreadout}
  \lean{PhyXMiniProblems.ProblemPhyXMini0191.acousticIntensityReadout}
  \uses{def:physics:phyx-mini-0191:phyxminiproblems-problemphyxmini0191-acousticintensityquantity}
  Read an acoustic intensity in a coherent choice of units
\end{definition}
\begin{proof}
  This is the declared model definition of acoustic Intensity Readout.
\end{proof}
\begin{definition}[distance In Meters]
  \label{def:physics:phyx-mini-0191:phyxminiproblems-problemphyxmini0191-distanceinmeters}
  \lean{PhyXMiniProblems.ProblemPhyXMini0191.distanceInMeters}
  \uses{def:physics:phyx-mini-0191:phyxminiproblems-problemphyxmini0191-lengthmagnitude, def:physics:phyx-mini-0191:phyxminiproblems-problemphyxmini0191-distancereadout}
  SI metre readout of a source-to-listener distance
\end{definition}
\begin{proof}
  This is the declared model definition of distance In Meters.
\end{proof}
\begin{definition}[acoustic Power In Watts]
  \label{def:physics:phyx-mini-0191:phyxminiproblems-problemphyxmini0191-acousticpowerinwatts}
  \lean{PhyXMiniProblems.ProblemPhyXMini0191.acousticPowerInWatts}
  \uses{def:physics:phyx-mini-0191:phyxminiproblems-problemphyxmini0191-acousticpowerquantity, def:physics:phyx-mini-0191:phyxminiproblems-problemphyxmini0191-acousticpowerreadout}
  SI watt readout of the emitted acoustic power
\end{definition}
\begin{proof}
  This is the declared model definition of acoustic Power In Watts.
\end{proof}
\begin{definition}[acoustic Intensity In Watts Per Square Meter]
  \label{def:physics:phyx-mini-0191:phyxminiproblems-problemphyxmini0191-acousticintensityinwattspersquaremeter}
  \lean{PhyXMiniProblems.ProblemPhyXMini0191.acousticIntensityInWattsPerSquareMeter}
  \uses{def:physics:phyx-mini-0191:phyxminiproblems-problemphyxmini0191-acousticintensityquantity, def:physics:phyx-mini-0191:phyxminiproblems-problemphyxmini0191-acousticintensityreadout}
  SI watt-per-square-metre readout of an acoustic intensity
\end{definition}
\begin{proof}
  This is the declared model definition of acoustic Intensity In Watts Per Square Meter.
\end{proof}
\begin{definition}[Listener Label]
  \label{def:physics:phyx-mini-0191:phyxminiproblems-problemphyxmini0191-listenerlabel}
  \lean{PhyXMiniProblems.ProblemPhyXMini0191.ListenerLabel}
  The two listener positions labeled in the supplied figure
\end{definition}
\begin{proof}
  This is the declared model definition of Listener Label.
\end{proof}
\begin{definition}[Figure Label]
  \label{def:physics:phyx-mini-0191:phyxminiproblems-problemphyxmini0191-figurelabel}
  \lean{PhyXMiniProblems.ProblemPhyXMini0191.FigureLabel}
  All labeled physical objects visible in the supplied figure
\end{definition}
\begin{proof}
  This is the declared model definition of Figure Label.
\end{proof}
\begin{definition}[Figure Role]
  \label{def:physics:phyx-mini-0191:phyxminiproblems-problemphyxmini0191-figurerole}
  \lean{PhyXMiniProblems.ProblemPhyXMini0191.FigureRole}
  Physical roles assigned to the three figure labels
\end{definition}
\begin{proof}
  This is the declared model definition of Figure Role.
\end{proof}
\begin{definition}[Wavefront Geometry]
  \label{def:physics:phyx-mini-0191:phyxminiproblems-problemphyxmini0191-wavefrontgeometry}
  \lean{PhyXMiniProblems.ProblemPhyXMini0191.WavefrontGeometry}
  Qualitative geometry of the blue sound-wave fronts in the image
\end{definition}
\begin{proof}
  This is the declared model definition of Wavefront Geometry.
\end{proof}
\begin{definition}[Acoustic Source Kind]
  \label{def:physics:phyx-mini-0191:phyxminiproblems-problemphyxmini0191-acousticsourcekind}
  \lean{PhyXMiniProblems.ProblemPhyXMini0191.AcousticSourceKind}
  Type of acoustic source used in the idealized model
\end{definition}
\begin{proof}
  This is the declared model definition of Acoustic Source Kind.
\end{proof}
\begin{definition}[Emission Regime]
  \label{def:physics:phyx-mini-0191:phyxminiproblems-problemphyxmini0191-emissionregime}
  \lean{PhyXMiniProblems.ProblemPhyXMini0191.EmissionRegime}
  Time dependence of the sound power emitted by the source
\end{definition}
\begin{proof}
  This is the declared model definition of Emission Regime.
\end{proof}
\begin{definition}[Acoustic Medium]
  \label{def:physics:phyx-mini-0191:phyxminiproblems-problemphyxmini0191-acousticmedium}
  \lean{PhyXMiniProblems.ProblemPhyXMini0191.AcousticMedium}
  The propagation medium shown in the outdoor scene
\end{definition}
\begin{proof}
  This is the declared model definition of Acoustic Medium.
\end{proof}
\begin{definition}[Bird Sound Figure]
  \label{def:physics:phyx-mini-0191:phyxminiproblems-problemphyxmini0191-birdsoundfigure}
  \lean{PhyXMiniProblems.ProblemPhyXMini0191.BirdSoundFigure}
  \uses{def:physics:phyx-mini-0191:phyxminiproblems-problemphyxmini0191-figurelabel, def:physics:phyx-mini-0191:phyxminiproblems-problemphyxmini0191-figurerole, def:physics:phyx-mini-0191:phyxminiproblems-problemphyxmini0191-wavefrontgeometry}
  Figure-specific qualitative labels and wavefront geometry
\end{definition}
\begin{proof}
  This is the declared model definition of Bird Sound Figure.
\end{proof}
\begin{definition}[Bird Point Source Sound Setup]
  \label{def:physics:phyx-mini-0191:phyxminiproblems-problemphyxmini0191-birdpointsourcesoundsetup}
  \lean{PhyXMiniProblems.ProblemPhyXMini0191.BirdPointSourceSoundSetup}
  \uses{def:physics:phyx-mini-0191:phyxminiproblems-problemphyxmini0191-lengthmagnitude, def:physics:phyx-mini-0191:phyxminiproblems-problemphyxmini0191-acousticpowerquantity, def:physics:phyx-mini-0191:phyxminiproblems-problemphyxmini0191-acousticintensityquantity, def:physics:phyx-mini-0191:phyxminiproblems-problemphyxmini0191-listenerlabel, def:physics:phyx-mini-0191:phyxminiproblems-problemphyxmini0191-acousticsourcekind, def:physics:phyx-mini-0191:phyxminiproblems-problemphyxmini0191-emissionregime, def:physics:phyx-mini-0191:phyxminiproblems-problemphyxmini0191-acousticmedium, def:physics:phyx-mini-0191:phyxminiproblems-problemphyxmini0191-birdsoundfigure}
  Physical quantities and observations for the two listener positions. There is one emitted-power field, expressing that the same source power is used at both positions. The two intensities remain independent until related by the inverse-square law; neither their ratio nor the requested decibel drop is stored in this structure
\end{definition}
\begin{proof}
  This is the declared model definition of Bird Point Source Sound Setup.
\end{proof}
\begin{definition}[Matches Point Source Scenario]
  \label{def:physics:phyx-mini-0191:phyxminiproblems-problemphyxmini0191-matchespointsourcescenario}
  \lean{PhyXMiniProblems.ProblemPhyXMini0191.MatchesPointSourceScenario}
  \uses{def:physics:phyx-mini-0191:phyxminiproblems-problemphyxmini0191-birdpointsourcesoundsetup}
  Textual physical model: an idealized bird is a point source with constant sound power, and sound propagates through ambient air
\end{definition}
\begin{proof}
  This is the declared model definition of Matches Point Source Scenario.
\end{proof}
\begin{definition}[Matches Supplied Figure]
  \label{def:physics:phyx-mini-0191:phyxminiproblems-problemphyxmini0191-matchessuppliedfigure}
  \lean{PhyXMiniProblems.ProblemPhyXMini0191.MatchesSuppliedFigure}
  \uses{def:physics:phyx-mini-0191:phyxminiproblems-problemphyxmini0191-distanceinmeters, def:physics:phyx-mini-0191:phyxminiproblems-problemphyxmini0191-birdpointsourcesoundsetup}
  Primary-image evidence. P₁ is the nearer listener, P₂ is the farther listener, and the blue curves represent concentric wavefronts centered on the bird. No numerical intensity ratio or decibel answer is asserted here
\end{definition}
\begin{proof}
  This is the declared model definition of Matches Supplied Figure.
\end{proof}
\begin{definition}[Moved To Twice The Distance]
  \label{def:physics:phyx-mini-0191:phyxminiproblems-problemphyxmini0191-movedtotwicethedistance}
  \lean{PhyXMiniProblems.ProblemPhyXMini0191.MovedToTwiceTheDistance}
  \uses{def:physics:phyx-mini-0191:phyxminiproblems-problemphyxmini0191-distancereadout, def:physics:phyx-mini-0191:phyxminiproblems-problemphyxmini0191-birdpointsourcesoundsetup}
  The stated move: the distance at P₂ is twice the distance at P₁ in every coherent choice of units. This premise contains no intensity or decibel conclusion
\end{definition}
\begin{proof}
  This is the declared model definition of Moved To Twice The Distance.
\end{proof}
\begin{definition}[Has Physical Acoustic Parameters]
  \label{def:physics:phyx-mini-0191:phyxminiproblems-problemphyxmini0191-hasphysicalacousticparameters}
  \lean{PhyXMiniProblems.ProblemPhyXMini0191.HasPhysicalAcousticParameters}
  \uses{def:physics:phyx-mini-0191:phyxminiproblems-problemphyxmini0191-distanceinmeters, def:physics:phyx-mini-0191:phyxminiproblems-problemphyxmini0191-acousticpowerinwatts, def:physics:phyx-mini-0191:phyxminiproblems-problemphyxmini0191-acousticintensityinwattspersquaremeter, def:physics:phyx-mini-0191:phyxminiproblems-problemphyxmini0191-listenerlabel, def:physics:phyx-mini-0191:phyxminiproblems-problemphyxmini0191-birdpointsourcesoundsetup}
  Positivity and nondegeneracy of the physical quantities in the model
\end{definition}
\begin{proof}
  This is the declared model definition of Has Physical Acoustic Parameters.
\end{proof}
\begin{definition}[Satisfies Isotropic Inverse Square Law]
  \label{def:physics:phyx-mini-0191:phyxminiproblems-problemphyxmini0191-satisfiesisotropicinversesquarelaw}
  \lean{PhyXMiniProblems.ProblemPhyXMini0191.SatisfiesIsotropicInverseSquareLaw}
  \uses{def:physics:phyx-mini-0191:phyxminiproblems-problemphyxmini0191-distancereadout, def:physics:phyx-mini-0191:phyxminiproblems-problemphyxmini0191-acousticpowerreadout, def:physics:phyx-mini-0191:phyxminiproblems-problemphyxmini0191-acousticintensityreadout, def:physics:phyx-mini-0191:phyxminiproblems-problemphyxmini0191-listenerlabel, def:physics:phyx-mini-0191:phyxminiproblems-problemphyxmini0191-birdpointsourcesoundsetup}
  The inverse-square intensity law for an isotropic point source. In every coherent unit system, the acoustic power through the sphere of radius r is 4 * pi * r\textasciicircum{}2 * I. This general law does not state the quarter-intensity or decibel conclusion for the doubled distance
\end{definition}
\begin{proof}
  This is the declared model definition of Satisfies Isotropic Inverse Square Law.
\end{proof}
\begin{definition}[sound Level Drop In Decibels]
  \label{def:physics:phyx-mini-0191:phyxminiproblems-problemphyxmini0191-soundleveldropindecibels}
  \lean{PhyXMiniProblems.ProblemPhyXMini0191.soundLevelDropInDecibels}
  \uses{def:physics:phyx-mini-0191:phyxminiproblems-problemphyxmini0191-acousticintensityinwattspersquaremeter, def:physics:phyx-mini-0191:phyxminiproblems-problemphyxmini0191-birdpointsourcesoundsetup}
  The drop in sound intensity level, in decibels, when moving from P₁ to P₂. This is the standard general conversion 10 log10 (I\_before / I\_after) applied to same-unit SI readouts; it contains no built-in value for the doubled-distance case
\end{definition}
\begin{proof}
  This is the declared model definition of sound Level Drop In Decibels.
\end{proof}
\begin{definition}[Answer Choice]
  \label{def:physics:phyx-mini-0191:phyxminiproblems-problemphyxmini0191-answerchoice}
  \lean{PhyXMiniProblems.ProblemPhyXMini0191.AnswerChoice}
  Labels of the four decibel choices printed in the problem
\end{definition}
\begin{proof}
  This is the declared model definition of Answer Choice.
\end{proof}
\begin{definition}[displayed Drop In Decibels]
  \label{def:physics:phyx-mini-0191:phyxminiproblems-problemphyxmini0191-displayeddropindecibels}
  \lean{PhyXMiniProblems.ProblemPhyXMini0191.displayedDropInDecibels}
  \uses{def:physics:phyx-mini-0191:phyxminiproblems-problemphyxmini0191-answerchoice}
  The decibel drop printed beside each answer label
\end{definition}
\begin{proof}
  This is the declared model definition of displayed Drop In Decibels.
\end{proof}
\begin{definition}[Matches Answer To Nearest Tenth Decibel]
  \label{def:physics:phyx-mini-0191:phyxminiproblems-problemphyxmini0191-matchesanswertonearesttenthdecibel}
  \lean{PhyXMiniProblems.ProblemPhyXMini0191.MatchesAnswerToNearestTenthDecibel}
  \uses{def:physics:phyx-mini-0191:phyxminiproblems-problemphyxmini0191-birdpointsourcesoundsetup, def:physics:phyx-mini-0191:phyxminiproblems-problemphyxmini0191-soundleveldropindecibels, def:physics:phyx-mini-0191:phyxminiproblems-problemphyxmini0191-answerchoice, def:physics:phyx-mini-0191:phyxminiproblems-problemphyxmini0191-displayeddropindecibels}
  Agreement with a displayed decibel value rounded to the nearest tenth. The half-tenth tolerance does not make choice A true by definition
\end{definition}
\begin{proof}
  This is the declared model definition of Matches Answer To Nearest Tenth Decibel.
\end{proof}
\begin{lemma}[intensity At P₂ eq one Quarter intensity At P₁]
  \label{lem:physics:phyx-mini-0191:phyxminiproblems-problemphyxmini0191-intensityatp-eq-onequarter-intensityatp}
  \lean{PhyXMiniProblems.ProblemPhyXMini0191.intensityAtP₂_eq_oneQuarter_intensityAtP₁}
  \uses{def:physics:phyx-mini-0191:phyxminiproblems-problemphyxmini0191-acousticintensityinwattspersquaremeter, def:physics:phyx-mini-0191:phyxminiproblems-problemphyxmini0191-birdpointsourcesoundsetup, def:physics:phyx-mini-0191:phyxminiproblems-problemphyxmini0191-movedtotwicethedistance, def:physics:phyx-mini-0191:phyxminiproblems-problemphyxmini0191-hasphysicalacousticparameters, def:physics:phyx-mini-0191:phyxminiproblems-problemphyxmini0191-satisfiesisotropicinversesquarelaw}
  Doubling the source distance makes the acoustic intensity one quarter as large. This is a derived intermediate result, not a model assumption
\end{lemma}
\begin{proof}
  The conclusion follows from the governing relations and figure readouts named in the statement of intensity At P₂ eq one Quarter intensity At P₁.
\end{proof}
% --- Archon declaration topology end ---
% --- Archon physics formalization source end ---
